\begin{lemma}[The crossing bonds of a native torus rectangle]%
    \label{lem:peps_torus_rectangle_boundary_card}
    \lean{TNLean.PEPS.card_horizontal_boundary_torusRectangle,
        TNLean.PEPS.card_vertical_boundary_torusRectangle,
        TNLean.PEPS.card_regionBoundaryEdge_torusRectangle,
        TNLean.PEPS.card_regionBoundaryEdge_torusSquare}
    \leanok
    \uses{def:peps_torus_geometry_translation,
        def:peps_torus_rectangular_grid}
    Suppose the native square torus has periods $W,H\geq3$. Let $R$
    be the coordinate rectangle
    \[
        [x_0,x_0+w)\times[y_0,y_0+h),
        \qquad w,h>0,
    \]
    with $x_0+w\leq W$, $y_0+h\leq H$, $w<W$, and $h<H$.
    Then there are $2h$ horizontal and $2w$ vertical crossing bonds,
    and hence
    \[
        |\partial R|=2w+2h.
    \]
    In particular a bounded square of side length $L$ has $4L$
    crossing bonds. The coordinate intervals may start at zero or
    end at their torus periods; the formula includes the cyclic seam
    bonds in these cases. These are the actual lattice bonds used in
    the rectangular boundary construction of
    \cite[Theorem~6.9]{Schuch2010PEPS}.
\end{lemma}

\begin{proof}\leanok
    For a cyclic coordinate interval $[s,s+\ell)$ shorter than its
    period $N$, membership changes between $a$ and $a+1$ at precisely
    two origins. Their coordinate values are $s+\ell-1$ and $s-1$,
    where the latter is read as $N-1$ if $s=0$.
    The two origins are distinct because $0<\ell<N$.
    A horizontal bond crosses the rectangle exactly when its first
    coordinate is one of these two origins and its second coordinate
    lies in the vertical interval. This gives $2h$ bonds. The vertical
    argument gives $2w$ bonds. On a torus with both periods at least
    three, right and up bonds enumerate the graph edges bijectively,
    so these two counts add without multiplicity.
\end{proof}

\begin{lemma}[A perimeter bound including full bands and empty rectangles]%
    \label{lem:peps_torus_rectangle_boundary_card_le}
    \lean{TNLean.PEPS.card_regionBoundaryEdge_torusRectangle_le}
    \leanok
    \uses{def:peps_torus_geometry_translation,
        def:peps_torus_rectangular_grid}
    On a native square torus with periods $W,H\geq3$, every bounded
    coordinate rectangle of side lengths $w,h\geq0$ satisfies
    $|\partial R|\leq2w+2h$. No strict upper bound on either side
    length is required; empty rectangles and full coordinate bands
    are included.
\end{lemma}

\begin{proof}\leanok
    \uses{lem:peps_torus_rectangle_boundary_card}
    A zero side length gives the empty rectangle. A side length equal
    to its torus period produces no crossing bond in that coordinate
    direction. In every other case, the horizontal crossing count is
    $2h$ and the vertical crossing count is $2w$. The two directional
    bounds add to the stated perimeter bound.
\end{proof}
